In an educational game where answering quiz questions is how the player
fights, a free-talking LLM non-player character can destroy the core loop by
handing out answers. We report an engineering case study of guarding such an
NPC. Stripping answer fields from its knowledge base does not prevent leakage:
66.8\% of the game's 500 quiz explanations state the correct option in prose.
We therefore adopt a context-graded threat model---discussing history is the
product, acting as an answer oracle is the failure---and guard request intent
rather than vocabulary. To evaluate, we pair every must-refuse case with a
must-answer case so that a mute NPC scores zero, and, because our development
cases were written while tuning the guardrails, we froze a blind held-out set
of 46 cases over three personas before running any model on it. Across a
cloud model, two local models and a LoRA-tuned 1.5B model (five rounds each),
development-set refusal holding of 82--100\% fell to 45--64\% on the held-out
set once leaks were judged by meaning and audited by hand. Every attack that
did not seek an answer was held in all 140 attempts; attacks asking for part
of an answer---elimination, a first character, an acrostic, confirmation of a
guess---mostly succeeded, and substring checks missed most of these leaks. We
distil five lessons on measurements that pointed the wrong way, and release
all code, cases and per-case run records.
